%%%%%%%%%%%%%%%%%%%%%%%%%%%%%%%%%%%%%%%%%
% Beamer Presentation - LaTeX Template
% Version 2.0 (March 8, 2022)
% Original Template: https://www.LaTeXTemplates.com
% Author: Vel (vel@latextemplates.com)
% License: CC BY-NC-SA 4.0

% Este modelo de apresentação foi 
% criado a partir do modelo de Giovanni Spadaro.
% Disponível em: https://github.com/Giovo17/presentation-template-unict-lm-data
%
% Adaptado por Lucas Amaral Taylor para criar uma versão especial 
% para os alunos de Matemática e Estatística da USP (IME-USP).
% Disponível em: https://github.com/lucasamtaylor01/IME-template
%%%%%%%%%%%%%%%%%%%%%%%%%%%%%%%%%%%%%%%%%

%----------------------------------------------------------------------------------------
% CLASSE DO DOCUMENTO E CONFIGURAÇÕES BÁSICAS
%----------------------------------------------------------------------------------------
\documentclass[
    11pt,               % Tamanho padrão da fonte
    % t,                % Alinhar verticalmente ao topo
    %aspectratio=169,   % Definir proporção 16:9
]{beamer}
\graphicspath{{img/}}         % Define o diretório das imagens

%----------------------------------------------------------------------------------------
% PACOTES NECESSÁRIOS
%----------------------------------------------------------------------------------------
\usepackage{
    booktabs,     % Melhora a aparência das linhas em tabelas
    palatino,     % Define Palatino como fonte principal
    subcaption    % Suporte para subfiguras
}
\usepackage[default]{opensans}  % Define Open Sans como fonte secundária
\input{config/code_langs}       % Importa configurações para highlight de código

%----------------------------------------------------------------------------------------
% CONFIGURAÇÃO DO TEMA
%----------------------------------------------------------------------------------------
% Tema Base
\usetheme{Boadilla}                          % Define o tema principal
\useinnertheme{circles}                      % Tema interno com círculos
\useoutertheme{miniframes}                   % Tema externo com miniframes
\setbeamertemplate{navigation symbols}{}     % Remove símbolos de navegação

% Cores Personalizadas
\definecolor{primaryColor}{RGB}{20,45,105}   % Cor primária - azul escuro
\definecolor{secondaryColor}{RGB}{0,100,160} % Cor secundária - azul médio

% Configurações de Cores
\setbeamercolor{structure}{fg=primaryColor}
\setbeamercolor{palette primary}{bg=primaryColor, fg=white}
\setbeamercolor{palette secondary}{bg=secondaryColor, fg=white}
\setbeamercolor{title}{bg=primaryColor, fg=white}

% Cores do Cabeçalho e Rodapé
\setbeamercolor{headline}{bg=secondaryColor, fg=white}
\setbeamercolor{section in head/foot}{bg=primaryColor, fg=white}
\setbeamercolor{subsection in head/foot}{bg=secondaryColor, fg=white}
\setbeamercolor{author in head/foot}{bg=primaryColor, fg=white}
\setbeamercolor{title in head/foot}{bg=secondaryColor, fg=white}
\setbeamercolor{date in head/foot}{bg=primaryColor, fg=white}
\setbeamercolor{page number in head/foot}{bg=primaryColor, fg=white}

%----------------------------------------------------------------------------------------
% BIBLIOGRAFIA
%----------------------------------------------------------------------------------------
\usepackage[style=alphabetic,backend=biber]{biblatex}
\addbibresource{bibliografia.bib}

%----------------------------------------------------------------------------------------
% INFORMAÇÕES DA APRESENTAÇÃO
%----------------------------------------------------------------------------------------
\title[Semanais]{Semanais}          % [Versão curta]{Versão completa}
\author[Gabriel Zocal]{Gabriel Zocal Santos}            % [Versão curta]{Nome completo}
\institute[Poli-USP]{Escola Politecnica \\ (Poli-USP)}
\date[2025]{10/25}

%----------------------------------------------------------------------------------------
% CAIXAS COLORIDAS (Tons equilibrados)
%----------------------------------------------------------------------------------------

\usepackage[most]{tcolorbox}
\usepackage{xcolor}

% Tons ajustados
\definecolor{Azul1}{HTML}{368FC5} % Teorema
\definecolor{Azul2}{HTML}{5EA9D6} % Definição
\definecolor{Azul3}{HTML}{8ABFE0} % Proposição
\definecolor{Azul4}{HTML}{B5D4EA} % Lema
\definecolor{Azul5}{HTML}{DBE9F4} % Corolário
\definecolor{Azul6}{HTML}{BFDCEC} % Observação
\definecolor{Azul7}{HTML}{A6CDE6} % Exemplo

% Caixa de Teorema
\newtcolorbox{CaixaTeo}[1]{colback=Azul1!10!white, colframe=Azul1, title={\textbf{Teorema:} #1}, coltitle=black}

% Caixa de Definição
\newtcolorbox{CaixaDef}[1]{colback=Azul2!10!white, colframe=Azul2, title={\textbf{Definição:} #1}, coltitle=black}

% Caixa de Proposição
\newtcolorbox{CaixaProp}[1]{colback=Azul3!10!white, colframe=Azul3, title={\textbf{Proposição:} #1}, coltitle=black}

% Caixa de Lema
\newtcolorbox{CaixaLema}[1]{colback=Azul4!10!white, colframe=Azul4, title={\textbf{Lema:} #1}, coltitle=black}

% Caixa de Corolário
\newtcolorbox{CaixaCoro}[1]{colback=Azul5!20!white, colframe=Azul5, title={\textbf{Corolário:} #1}, coltitle=black}

% Caixa de Observação
\newtcolorbox{CaixaObs}[1]{colback=Azul6!10!white, colframe=Azul6, title={\textbf{Observação:} #1}, coltitle=black}

% Caixa de Exemplo
\newtcolorbox{CaixaExem}[1]{colback=Azul7!10!white, colframe=Azul7, title={\textbf{Exemplo:} #1}, coltitle=black}

%----------------------------------------------------------------------------------------
% INÍCIO DO DOCUMENTO
%----------------------------------------------------------------------------------------
\begin{document}

% Slide de título com logo
\begin{frame}
    \begin{figure}
        \includegraphics[width=0.4\linewidth]{img/logo_POLI2.png}
    \end{figure}
    \titlepage
\end{frame}

% Sumário
\begin{frame}
    \frametitle{Estrutura da apresentação}
    \tableofcontents
    
\end{frame}

% Inclusão das seções
\section{Texto} % Seções são adicionadas para organizar sua apresentação em blocos discretos, todas as seções e subseções são automaticamente exibidas no índice como uma visão geral da apresentação, mas NÃO são exibidas como slides separados.

%----------------------------------------------------------------------------------------

% \begin{frame}
% 	\frametitle{Texto corrido}
%     Lorem ipsum dolor sit amet, consectetur adipiscing elit. Nullam ipsum velit, cursus quis ligula eu, malesuada aliquet massa. Quisque non convallis felis, a auctor eros. Etiam sit amet turpis a sapien pulvinar malesuada quis quis nisi. Quisque scelerisque volutpat ligula vel mollis. Nam sit amet tristique erat, sit amet cursus mi. 
% \end{frame}

%----------------------------------------------------------------------------------------

% \begin{frame}
% 	\frametitle{Texto em tópicos numerados}
%      Lorem ipsum dolor sit amet, consectetur adipiscing elit:
%     \begin{enumerate}
%         \item Lorem ipsum dolor sit amet.
%         \item Lorem ipsum dolor sit amet.
%     \end{enumerate}
	
% \end{frame}

%----------------------------------------------------------------------------------------

\begin{frame}
	\frametitle{Conceitos novos}
    Familiarização com novos conceitos
    \begin{itemize}
        \item Perfect security
        \item Limitações one time pad
        \item Semantic secure 
        \item Feistel NT,DES, 3DES
        \item PRF, Block cipher 
    \end{itemize}
	
\end{frame}


\begin{frame}
    \frametitle{FHE}
    Novos topicos de Matemática
    \begin{itemize}
        \item Group Theory
        \item Polynomial Ring 
    \end{itemize}
\end{frame}
%% \section{Equações e imagens} % Seções são adicionadas para organizar sua apresentação em blocos discretos, todas as seções e subseções são automaticamente exibidas no índice como uma visão geral da apresentação, mas NÃO são exibidas como slides separados.

% %----------------------------------------------------------------------------------------

% \begin{frame}{Uma imagem}
%    \begin{figure}[h]
%        \centering
%        \includegraphics[width=0.7\textwidth]{img/logo_IME.png}
%        \caption{Legenda da imagem}
%        \label{fig:label_da_imagem}
%    \end{figure}
% \end{frame}

% %----------------------------------------------------------------------------------------

% \begin{frame}{Duas imagens}
%    \begin{figure}
%        \centering
%        \begin{subfigure}[b]{0.45\textwidth}
%            \centering
%            \includegraphics[width=\textwidth]{img/logo_IME.png}
%            \caption{Legenda 1}
%            \label{fig:img1}
%        \end{subfigure}
%        \hfill
%        \begin{subfigure}[b]{0.45\textwidth}
%            \centering
%            \includegraphics[width=\textwidth]{img/logo_IME.png}
%            \caption{Legenda 2}
%            \label{fig:img2}
%        \end{subfigure}
%    \end{figure}
% \end{frame}

%----------------------------------------------------------------------------------------
% \begin{frame}{Equações}
%     Equações de Navier-Stokes Forma expandida (3D):
%     \footnotesize
%         \begin{align*}
%             \rho\left(\frac{\partial u}{\partial t} + u\frac{\partial u}{\partial x} + v\frac{\partial u}{\partial y} + w\frac{\partial u}{\partial z}\right) &= -\frac{\partial p}{\partial x} + \mu\left(\frac{\partial^2 u}{\partial x^2} + \frac{\partial^2 u}{\partial y^2} + \frac{\partial^2 u}{\partial z^2}\right) + f_x \\[0.3cm]
%             \rho\left(\frac{\partial v}{\partial t} + u\frac{\partial v}{\partial x} + v\frac{\partial v}{\partial y} + w\frac{\partial v}{\partial z}\right) &= -\frac{\partial p}{\partial y} + \mu\left(\frac{\partial^2 v}{\partial x^2} + \frac{\partial^2 v}{\partial y^2} + \frac{\partial^2 v}{\partial z^2}\right) + f_y \\[0.3cm]
%             \rho\left(\frac{\partial w}{\partial t} + u\frac{\partial w}{\partial x} + v\frac{\partial w}{\partial y} + w\frac{\partial w}{\partial z}\right) &= -\frac{\partial p}{\partial z} + \mu\left(\frac{\partial^2 w}{\partial x^2} + \frac{\partial^2 w}{\partial y^2} + \frac{\partial^2 w}{\partial z^2}\right) + f_z
%         \end{align*}
        
%     onde $\mathbf{v} = (u,v,w)$ é o campo de velocidade, $p$ é a pressão, $\rho$ é a densidade, $\mu$ é a viscosidade dinâmica e $\mathbf{f}$ representa forças externas.
% \end{frame}
% % \section{Código} % Seções são adicionadas para organizar sua apresentação em blocos discretos, todas as seções e subseções são automaticamente exibidas no índice como uma visão geral da apresentação, mas NÃO são exibidas como slides separados.

% %----------------------------------------------------------------------------------------
% \begin{frame}[fragile]
%     \frametitle{Python}
    
%     \begin{python}
% def calcular_dobro(x):
%     """Retorna o dobro do número"""
%     return 2 * x

% # Testando a função
% numero = 5
% resultado = calcular_dobro(numero)
% print(f"O dobro de {numero} é {resultado}")
%     \end{python}
% \end{frame}

% %----------------------------------------------------------------------------------------
% \begin{frame}[fragile]
%     \frametitle{C}
    
%     \begin{clang}
% #include <stdio.h>

% int main() {
%     int numero = 5;
%     int dobro = 2 * numero;
    
%     printf("O dobro de %d eh %d\n", numero, dobro);
%     return 0;
% }
%     \end{clang}
% \end{frame}

% %----------------------------------------------------------------------------------------
% \begin{frame}[fragile]
%     \frametitle{C++}
    
%     \begin{cpp}
% #include <iostream>
% using namespace std;

% int main() {
%     int numero = 5;
%     int dobro = 2 * numero;
    
%     cout << "O dobro de " << numero;
%     cout << " eh " << dobro << endl;
%     return 0;
% }
%     \end{cpp}
% \end{frame}

% %----------------------------------------------------------------------------------------
% \begin{frame}[fragile]
%     \frametitle{R}
    
%     \begin{rlang}
% # Função para calcular o dobro
% calcular_dobro <- function(x) {
%   return(2 * x)
% }

% # Testando a função
% numero <- 5
% resultado <- calcular_dobro(numero)
% print(paste("O dobro de", numero, "é", resultado))
%     \end{rlang}
% \end{frame}

% %----------------------------------------------------------------------------------------

% \begin{frame}[fragile]
%     \frametitle{Java}
    
%     \begin{java}
% public class Exemplo {
%     public static void main(String[] args) {
%         int numero = 5;
%         int dobro = 2 * numero;
        
%         System.out.println("O dobro de " + numero +
%                          " eh " + dobro);
%     }
% }
%     \end{java}
% \end{frame}

% % \section{Caixas coloridas}

% \begin{frame}{Caixas: Teorema e Definição}
% 	\begin{CaixaTeo}{Teorema Fundamental da Aritmética}
% 		Todo inteiro $n > 1$ pode ser escrito de forma única como um produto de primos, a menos da ordem dos fatores.
% 	\end{CaixaTeo}
	
% 	\begin{CaixaDef}{Número Primo}
% 		Um número natural $p > 1$ é primo se seus únicos divisores positivos são $1$ e $p$.
% 	\end{CaixaDef}
% \end{frame}

% \begin{frame}{Caixas: Proposição e Lema}
% 	\begin{CaixaProp}{Divisibilidade}
% 		Se $a \mid b$ e $b \mid c$, então $a \mid c$.
% 	\end{CaixaProp}
	
% 	\begin{CaixaLema}{Lema de Euclides}
% 		Se um primo $p$ divide o produto $ab$, então $p$ divide $a$ ou $p$ divide $b$.
% 	\end{CaixaLema}
% \end{frame}

% \begin{frame}{Caixas: Corolário e Observação}
% 	\begin{CaixaCoro}{Infinidade de Primos}
% 		Existem infinitos números primos.
% 	\end{CaixaCoro}
	
% 	\begin{CaixaObs}{Máximo Divisor Comum}
% 		O máximo divisor comum de dois inteiros pode ser calculado pelo Algoritmo de Euclides.
% 	\end{CaixaObs}
% \end{frame}

% \begin{frame}{Caixas: Exemplo e Nota}
% 	\begin{CaixaExem}{MDC com Algoritmo de Euclides}
% 		Para calcular $\gcd(48,18)$:
% 		\begin{equation*}
% 			48 = 2 \cdot 18 + 12,\quad 18 = 1 \cdot 12 + 6,\quad 12 = 2 \cdot 6 + 0.
% 		\end{equation*}
% 		Logo, $\gcd(48,18) = 6$.
% 	\end{CaixaExem}
	
% 	\begin{CaixaObs}{Fato curioso}
% 		O número $26$ é o único número natural que está entre um cubo perfeito e um quadrado perfeito.
% 	\end{CaixaObs}
	
% \end{frame}
\end {frame}
%% \section{Conclusão} % Seções são adicionadas para organizar sua apresentação em blocos discretos, todas as seções e subseções são automaticamente exibidas no índice como uma visão geral da apresentação, mas NÃO são exibidas como slides separados.

% \begin{frame}{Referências}
%     \nocite{*}
%     \printbibliography[heading=none]
% \end{frame}

% Slide final
\begin{frame}
    \begin{center}
        {\Huge Fim da apresentação!}
    \end{center}
\end{frame}

\end{document}


